\BeleskeID{09_3-s061}
% element: 09_3-s061:e05
% slide-id: 09_3-s061
N tranzistor ima velik napon između gejta i sorsa. Sors je na ulaznoj strani 1, a struja kroz $R_L$ teče iz mase ka negativnom ulazu u smeru označenom na šemi. Drejn ostaje na negativnom potencijalu; zbog malog napona drejn–sors N tranzistor radi u omskoj oblasti. Za opterećenje mnogo veće od njegove dinamičke otpornosti izlaz je približno $-12\,\mathrm V$, bez degradacije nivoa kao kod samog prolaznog N tranzistora koji prenosi visoku jedinicu.

Pod pretpostavkom $V_{Tn}\approx-V_{Tp}<3\,\mathrm V$, i P tranzistor ima uslove za provođenje. Njegov sors je na strani 2 i važi $V_{GS}<V_{Tp}$. N tranzistor održava malu razliku potencijala između krajeva, pa P tranzistor takođe radi u omskoj oblasti. Zbog manje razlike sors–gejt njegova otpornost je veća.
